% !TeX root=../main.tex
% در این فایل، عنوان پایان‌نامه، مشخصات خود، متن تقدیمی‌، ستایش، سپاس‌گزاری و چکیده پایان‌نامه را به فارسی، وارد کنید.
% توجه داشته باشید که جدول حاوی مشخصات پروژه/پایان‌نامه/رساله و همچنین، مشخصات داخل آن، به طور خودکار، درج می‌شود.
%%%%%%%%%%%%%%%%%%%%%%%%%%%%%%%%%%%%
% دانشگاه خود را وارد کنید
\university{دانشگاه خواجه نصیرالدین طوسی}
% پردیس دانشگاهی خود را اگر نیاز است وارد کنید (مثال: فنی، علوم پایه، علوم انسانی و ...)
\college{دانشکدهٔ مهندسی کامپیوتر}
% دانشکده، آموزشکده و یا پژوهشکده  خود را وارد کنید
\faculty{دانشکدهٔ مهندسی کامپیوتر}
% گروه آموزشی خود را وارد کنید (در صورت نیاز)
\department{گروه هوش مصنوعی}
% رشته تحصیلی خود را وارد کنید
\subject{مهندسی کامپیوتر}
% در صورت داشتن گرایش، خط زیر را از حالت کامت خارج نموده و گرایش خود را وارد کنید
% \field{گرابش}
% عنوان پایان‌نامه را وارد کنید
\title{تخمین عمق بدون نظارت با استفاده از پارامترهای کالیبراسیون داخلی به عنوان ورودی}
% نام استاد(ان) راهنما را وارد کنید
\firstsupervisor{دکتر بهروز نصیحت کن}
\firstsupervisorrank{استادیار}
%\secondsupervisor{}
%\secondsupervisorrank{}
% نام استاد(دان) مشاور را وارد کنید. چنانچه استاد مشاور ندارید، دستورات پایین را غیرفعال کنید.
%\firstadvisor{دکتر مشاور اول}
%\firstadvisorrank{استادیار}
%\secondadvisor{دکتر مشاور دوم}
% نام داوران داخلی و خارجی خود را وارد نمایید.
\internaljudge{دکتر حمید ابریشمی مقدم}
\internaljudgerank{استاد}
\externaljudge{دکتر محمدرضا محمدی}
\externaljudgerank{استادیار}
% نام دانشگاه داور خارجی در فرم چاپی این قالب نمایش داده نمی‌شود.
% نام نماینده کمیته تحصیلات تکمیلی در دانشکده \ گروه
% مشخصات نمایندهٔ تحصیلات تکمیلی پس از تعیین رسمی افزوده می‌شود.
% نام دانشجو را وارد کنید
\name{محمدامین}
% نام خانوادگی دانشجو را وارد کنید
\surname{قاسم‌زاده}
% شماره دانشجویی دانشجو را وارد کنید
\studentID{40209644}
% تاریخ پایان‌نامه را وارد کنید
\thesisdate{تابستان ۱۴۰۵}
% به صورت پیش‌فرض برای پایان‌نامه‌های کارشناسی تا دکترا به ترتیب از عبارات «پروژه»، «پایان‌نامه» و «رساله» استفاده می‌شود؛ اگر  نمی‌پسندید هر عنوانی را که مایلید در دستور زیر قرار داده و آنرا از حالت توضیح خارج کنید.
%\projectLabel{پایان‌نامه}

% به صورت پیش‌فرض برای عناوین مقاطع تحصیلی کارشناسی تا دکترا به ترتیب از عبارت «کارشناسی»، «کارشناسی ارشد» و «دکتری» استفاده می‌شود؛ اگر نمی‌پسندید هر عنوانی را که مایلید در دستور زیر قرار داده و آنرا از حالت توضیح خارج کنید.
%\degree{}
%%%%%%%%%%%%%%%%%%%%%%%%%%%%%%%%%%%%%%%%%%%%%%%%%%%%
%% پایان‌نامه خود را تقدیم کنید! %%
\dedication
{
{\Large تقدیم به: مادر و پدر عزیزم}\\
\begin{flushleft}{
	\huge
به آنان که با علم خود زندگی آزاد می‌سازند\\
	\vspace{7mm}
}
\end{flushleft}
}
%% متن قدردانی %%
%% این متن را به سلیقه‌ی خود تعییر دهید
\acknowledgement{
اکنون که به یاری پروردگار و یاری و راهنمایی اساتید بزرگ موفق به پایان این رساله شده‌ام وظیفه خود دانشته که نهایت سپاسگزاری را از تمامی عزیزانی که در این راه به من کمک کرده‌اند را به عمل آورم:
در آغاز از استاد بزرگ و دانشمند جناب آقادکتر نصیحت‌کن که راهنمایی این پایان‌نامه را به عهده داشته‌اند کمال تشکر را دارم.
از داوران گرامی دکتر حمید ابریشمی مقدم و دکتر محمدرضا محمدی که زحمت داوری و تصحیح این پایان‌نامه را به عهده داشتند کمال سپاس را دارم.
خالصانه از تمامی اساتید و معلمان و مدرسانی که در مقاطع مختلف تحصیلی به من علم آموخته و مرا از سرچشمه دانایی سیراب کرده‌اند متشکرم.
از کلیه هم دانشگاهیان و همراهان عزیز، دوستان خوبم نهایت سپاس را دارم.

و در پایان این پایان‌نامه را تقدیم می‌کنم به مادر و پدر مهربانم که با حضور و همراهی شان در دوران تحصیل، همیشه راه را به من نشان داده و مرا در این راه استوار و ثابت قدم نموده است.
}
%%%%%%%%%%%%%%%%%%%%%%%%%%%%%%%%%%%%
%چکیده پایان‌نامه را وارد کنید
\faabstract{
تخمین عمق تک‌دوربینه ذاتاً مسئله‌ای مبهم است، زیرا یک تصویر منفرد مقیاس متریک صحنه را به‌صورت یکتا تعیین نمی‌کند. در این پایان‌نامه بررسی می‌شود که آیا ورود صریح پارامترهای درونی دوربین می‌تواند این ابهام را در یک مدل بنیادین تخمین عمق کاهش دهد. برای این منظور، نسخه‌ای آگاه از دوربین از مدل \lr{Depth Anything V2} توسعه داده شده است که در آن یک توصیف‌گر نه‌بعدی نرمال‌شده از ماتریس مشخصهٔ دوربین، به توکنی دروازه‌بانی‌شده تبدیل و به انکودر \lr{DINOv2} تزریق می‌شود. روش تجربی طی مجموعه‌ای از آزمایش‌های کنترل‌شده تکامل یافت و در این مسیر، نشت داده در تفکیک \lr{NYUv2}، ناهم‌خوانی خاموش میان فضای خروجی دیسپاریتی و هدف عمق مستقیم، به‌روزرسانی‌نشدن $K$ هنگام برش تصویر و ناهم‌خوانی خاموش سر خروجی مدل متریک در ارزیابی شناسایی و اصلاح شد.

تمام مقایسه‌های نهایی بر مجموعهٔ آزمون بدون نشت ۶۵۴ تصویری \lr{Eigen} و میان دو بازوی آموزشی همسان، \proposedmodel{} و \controlmodel{}، انجام شده‌اند. در پروتکل نسبی، \proposedmodel{}، نسخهٔ \lr{ViT-S}، از \publishedbasemodel{} و \controlmodel{} بهتر است؛ با این حال، در نسخهٔ \lr{ViT-L} سهم خالص $K$ معنادار نیست، زیرا تراز مقیاس و جابجایی به‌ازای هر تصویر، مهم‌ترین اطلاعات حاصل از فاصلهٔ کانونی را حذف می‌کند. در ارزیابی متریک بدون تراز، \proposedmodel{} به‌طور پایدار از \controlmodel{} بهتر عمل می‌کند و با افزایش تنوع میدان دید، اندازهٔ اثر نیز افزایش می‌یابد. سپس افزونه‌سازی هندسی «لرزش کانونی» پیشنهاد می‌شود که بازیابی مقیاس را تنها از روی پیکسل‌ها ناممکن می‌سازد. در این تنظیم، \proposedmodel{} به AbsRel برابر \fadecimal{۰٫۰۸۴۹} در برابر \fadecimal{۰٫۱۰۳۰} برای \controlmodel{} می‌رسد و در همهٔ معیارهای گزارش‌شده با $p\leq10^{-24}$ برتری دارد.

آزمون‌های سازوکاری و علّی نشان می‌دهند که این بهبود واقعاً ناشی از استفاده از ماتریس مشخصه است. شیب لگاریتمی پاسخ عمق پیش‌بینی‌شده به مقیاس فاصلهٔ کانونی، پس از لرزش کانونی از \fadecimal{۰٫۰۶} به \fadecimal{۰٫۹۳} تا \fadecimal{۰٫۹۷} می‌رسد که به مقدار ایده‌آل یک در مدل روزنه‌ای نزدیک است. حذف $K$ خطای AbsRel را ۵۹ درصد افزایش می‌دهد و تحریف عمدی فاصلهٔ کانونی، خطایی هم‌اندازه با پیش‌بینی هندسهٔ تصویربرداری ایجاد می‌کند. بنابراین، پارامترهای درونی دوربین می‌توانند سیگنالی معنادار و از نظر علّی ضروری برای بازیابی مقیاس فراهم کنند؛ هرچند این اتکا، حساسیت به دقت کالیبراسیون را نیز افزایش می‌دهد و تعمیم آن به داده‌های واقعی چنددوربینه و محیط‌های بیرونی نیازمند بررسی‌های آینده است.

}
% کلمات کلیدی پایان‌نامه را وارد کنید
\keywords{
تخمین عمق تک‌دوربینه،
یادگیری عمیق بدون نظارت،
عمق متریک،
پارامترهای درونی دوربین،
هندسهٔ تصویربرداری،
Depth Anything V2
}
% انتهای وارد کردن فیلد‌ها
%%%%%%%%%%%%%%%%%%%%%%%%%%%%%%%%%%%%%%%%%%%%%%%%%%%%%%
